\section{Objetivo experimental (1.1)}

El experimento compara métodos de generación de kernels en Triton DSL usando SLM locales y modelos de frontera, con un pipeline de verificación antes de compilar en GPU. El sistema del equipo, Triton Heal en modo dual, genera con Llama-3.1-8B y, si el código no pasa las validaciones, intenta repararlo con DeepSeek-Coder-V2 (hasta dos reintentos).

Se contrastan seis configuraciones: solo Llama-8B, solo DeepSeek-R1-Distill-Qwen-14B, solo DeepSeek-Coder-V2, Claude 3.5 Sonnet, GPT-4o y el dual propuesto. Las preguntas centrales son si el dual mejora la tasa de compilación frente a GPT-4o y frente a Llama solo, y si el filtrado por verificación (MSR, FNR, latencia) aporta frente a usar un único SLM.

\section{Hipótesis (1.2)}

Todas las pruebas usan $\alpha = 0.05$ antes de corrección; los $p$-valores ajustados siguen Holm-Bonferroni (sección 1.8).

\begin{table}[ht]
\centering
\caption{Hipótesis por contraste y métrica.}
\label{tab:hipotesis}
\footnotesize
\begin{adjustbox}{max width=\linewidth,center}
\begin{tabularx}{\linewidth}{@{}lXX@{}}
\toprule
Contraste & $H_0$ & $H_1$ \\
\midrule
C4: dual vs GPT-4o (compilación) & Misma tasa de compilación & Dual con mayor tasa \\
C7: dual vs Llama (compilación) & Misma tasa & Dual con mayor tasa \\
C1: Claude vs Llama (MSR) & Mismo MSR & Claude con mayor MSR \\
C2: GPT vs Coder-V2 (MSR) & Mismo MSR & GPT con mayor MSR \\
C3: dual vs Llama (FNR) & Mismo FNR & Dual con menor FNR \\
C5: latencia Claude vs Llama & Misma mediana & Frontera más lenta \\
C8: speedup proxy dual vs GPT & Igual & Dual mayor \\
\bottomrule
\end{tabularx}
\end{adjustbox}
\end{table}

\section{Variables experimentales (1.3)}

\subsection{Variables independientes}

La configuración del sistema (seis modos), el modelo concreto (SLM local vs frontera), el prompt (\texttt{generator\_v1.txt} o \texttt{generator\_repair\_v1.txt} en reparación), la temperatura ($T = 0$) y el modo de ejecución (\texttt{VERIFIER\_MODE}, típicamente Ollama en nuestro entorno).

\subsection{Variables dependientes}

En generación: si el código compila bajo nuestro checker (\texttt{compile\_ok}), la tasa de compilación en porcentaje, el speedup proxy y un indicador de eficiencia operativa del pipeline. En verificación del benchmark fijo: MSR, FNR, F1, porcentaje de JSON válido y latencia en milisegundos.

\subsection{Variables controladas}

Mismo benchmark de 70 kernels y el mismo conjunto de tareas de generación para todas las configuraciones; hardware Apple M4 con 24\,GB; \texttt{seed} 42; \texttt{timeout} 120\,s por inferencia. El diseño es within-subjects: cada tarea o kernel se evalúa en todos los métodos.

\section{Diseño experimental (1.4)}

Cada tarea de generación se asigna a las seis configuraciones (en la corrida archivada, R1-14B se omitió por timeouts). El plan de diseño contempla 20 tareas con tres réplicas (360 ejecuciones de generación) y 70 kernels con tres réplicas en verificación (1260 evaluaciones). Se eligió within-subjects para poder aplicar McNemar y Wilcoxon pareados y reducir la varianza entre tareas de distinta dificultad.

\section{Métricas de evaluación (1.5)}

\subsection{Tasa de compilación}

Proporción de salidas con \texttt{compile\_ok}. En esta entrega el criterio exige pasar tres niveles en Mac M4: sintaxis Python (AST), reglas léxicas del checker del equipo, y presencia de \texttt{@triton.jit} con operaciones \texttt{tl.load}/\texttt{tl.store}. No equivale a compilar en CUDA ni a ejecutar en GPU.

\subsection{Speedup proxy}

Se define $s = t_{\mathrm{baseline}} / t_{\mathrm{pipeline}}$, donde $t_{\mathrm{baseline}}$ es el tiempo del checker sobre el kernel de referencia del benchmark y $t_{\mathrm{pipeline}}$ suma la latencia de generación más la del checker sobre el código generado. Sirve como indicador operativo en M4; el informe no reporta speedup medido en GPU NVIDIA.

\subsection{Eficiencia operativa (proxy)}

Relación entre compilaciones exitosas y tiempo de generación. No sustituye métricas de ocupación de SM o ancho de banda de memoria.

\subsection{Verificación (MSR, FNR, F1, JSON válido, latencia)}

MSR es la fracción de kernels \texttt{unsafe} correctamente marcados como inseguros; FNR es la fracción de \texttt{unsafe} pasados por alto. F1 resume precisión y recall en la clasificación safe/unsafe. El porcentaje de JSON válido indica si el verificador devolvió una respuesta parseable.

\section{Tamaño del efecto esperado (1.6)}

En contrastes pareados de compilación y MSR se anticipa un efecto mediano ($\phi \approx 0.4$ en McNemar). Además de la significación estadística, se exige una diferencia práctica de al menos 5 puntos porcentuales en compilación o MSR, y alrededor de 8 pp de reducción de FNR para considerar relevante la mejora del dual.

\section{Análisis de poder (1.7)}

Nivel de significación $\alpha = 0.05$ y poder objetivo $1 - \beta = 0.80$. La muestra de generación tiene 20 tareas; en verificación hay 31 kernels \texttt{unsafe}, cantidad suficiente para detectar efectos medianos en McNemar según el diseño del curso.

\section{Plan estadístico (1.8)}

\begin{table}[ht]
\centering
\caption{Contrastes confirmatorios y prueba asociada.}
\label{tab:plan_stats}
\footnotesize
\begin{tabularx}{\linewidth}{@{}llX@{}}
\toprule
ID & Prueba & Comparación \\
\midrule
C4 & McNemar & Compilación: dual vs GPT-4o \\
C7 & McNemar & Compilación: dual vs Llama \\
C8 & Wilcoxon & Speedup proxy: dual vs GPT \\
C1--C3, C5 & McNemar / Wilcoxon & Verificación en benchmark \\
C6 & Friedman + Holm & F1 entre configuraciones \\
\bottomrule
\end{tabularx}
\end{table}

Se rechaza $H_0$ cuando $p_{\mathrm{Holm}} < 0.05$ y se cumple el umbral de relevancia práctica acordado. Para proporciones se reportan intervalos de confianza Wilson al 95\% (MSR en la tabla extendida de la Parte II). La familia de contrastes se corrige con Holm-Bonferroni.

\subsection{Desviaciones respecto al plan inicial}

Durante la ejecución se corrigió la taxonomía de DeepSeek (Coder-V2 como frontera, R1-14B como SLM local). Claude y GPT corrieron con verificador heurístico al no disponer de API en el entorno de prueba. La generación archivada usa 10 tareas y excluye \texttt{solo\_deepseek} por timeout; la verificación quedó con una réplica por kernel (\texttt{run\_id = 0}), aunque \texttt{experiment\_config.json} ya deja configuradas tres réplicas para corridas futuras.

\section{Reproducibilidad experimental (1.9)}

\begin{table}[ht]
\centering
\caption{Entorno y artefactos.}
\label{tab:repro}
\footnotesize
\begin{adjustbox}{max width=\linewidth,center}
\begin{tabularx}{\linewidth}{@{}lX@{}}
\toprule
Componente & Especificación \\
\midrule
Repositorio & \url{https://github.com/G4ndhiV/TRITONHEAL} (commit \texttt{c05f82a}) \\
Hardware & macOS 15.x; Apple M4; 24\,GB memoria unificada \\
Ollama & 0.24.x (\texttt{llama3.1:8b}, \texttt{deepseek-r1:14b}, \texttt{deepseek-coder-v2}) \\
Modelos API & Claude 3.5 Sonnet; GPT-4o (heurístico si no hay clave) \\
Python & 3.10+; dependencias en \texttt{requirements.txt} \\
Benchmark & 70 kernels; \texttt{generation\_tasks.jsonl} \\
Prompts & \texttt{prompts/generator\_v1.txt}, \texttt{generator\_repair\_v1.txt} \\
Reproducción & \texttt{make benchmark generation experiments stats-all figures pdf} \\
\bottomrule
\end{tabularx}}
\end{adjustbox}
\end{table}

\section{Amenazas a la validez (pre-registro, 1.10)}

\begin{table}[ht]
\centering
\caption{Amenazas anticipadas y mitigación.}
\label{tab:amenazas_pre}
\footnotesize
\begin{adjustbox}{max width=\linewidth,center}
\begin{tabularx}{\linewidth}{@{}lXX@{}}
\toprule
Tipo & Riesgo & Mitigación \\
\midrule
Interna & Mezcla API real y heurístico & Tabla planeada vs ejecutado en apéndice \\
Interna & Timeouts en modelos grandes & Fijar timeout; reportar JSON válido \\
Externa & Solo Mac M4 (Metal) & No extrapolar a clusters NVIDIA \\
Externa & Benchmark sintético & Etiquetar categorías de error \\
Constructo & Compilación $\neq$ corrección GPU & Definir niveles L1--L3 \\
Constructo & Speedup proxy & Declarar limitación en figuras \\
\bottomrule
\end{tabularx}
\end{adjustbox}
\end{table}
